\subsubsection{Determinism}
\label{sec:semantics:grsync:theorems:determinism}

A deterministic component produces a unique output history from a given input history. In other
words, given the same input, the component will always produce the same output.

This property is formally modeled by \Cref{thm:semantics:grsync:theorems:determinism}: the
denotational semantics of a component, \ie its associated stream function, ensures that two output
streams mapped to one input stream are identical.
%
Proving the equality of such infinite streams requires a coinductive argument, which we structure as
an induction on the length of the stream prefixes.

To this end, we introduce an intermediate equivalence relation, denoted by $\equiv_n$ and defined in
\Cref{def:semantics:grsync:theorems:determinism:equiv_upto_n}, which captures the equality of two
streams up to their first $n$ elements.

\begin{definition}[Equivalence up-to-$n$]
    \label{def:semantics:grsync:theorems:determinism:equiv_upto_n}
    Let $n \in \N$.
    The relation $\equiv_n$ denotes the equality of two streams up-to-$n$ elements and is defined
    by induction on $n$:
    \begin{align}
        \equiv_0 & = \Stream{\ValBot} \times \Stream{\ValBot} \label{def:semantics:grsync:theorems:determinism:equiv0} \\
        \equiv_n & = \{ (v_{\bot1} \step s_1, v_{\bot2} \step s_2) \mid s_1 \equiv_{n-1} s_2,
        v_{\bot1}, v_{\bot2} \in V + \bot, v_{\bot1} = v_{\bot2} \} \label{def:semantics:grsync:theorems:determinism:equivn}
    \end{align}
\end{definition}

\begin{lemma}
    Let $n \in \N$.
    The relation $\equiv_n$ is an equivalence relation.
\end{lemma}
\begin{skippableproof}
    We reason by induction on $n$.
    \begin{itemize}
        \item Reflexive: \Cref{def:semantics:grsync:theorems:determinism:equiv0} implies for all stream $s$ we
              have $s \equiv_0 s$.

              Let $n \in \N$ and suppose that for all stream $s$ we have $s \equiv_n s$.
              Let $s = v_\bot \step s'$ be a stream. The induction hypothesis implies that
              $s' \equiv_n s'$, and \Cref{def:semantics:grsync:theorems:determinism:equivn} ensures
              $v_\bot \step s' \equiv_{n+1} v_\bot \step s'$, so that we have $s \equiv_{n+1} s$.

        \item Symmetrical: \Cref{def:semantics:grsync:theorems:determinism:equiv0} implies for all streams
              $s_1$ and $s_2$ we have $s_1 \equiv_0 s_2$ and $s_2 \equiv_0 s_1$.

              Let $n \in \N$ and suppose that for all streams $s_1$ and $s_2$ such that
              $s_1 \equiv_n s_2$ we have $s_2 \equiv_n s_1$.
              Let $s_1 = v_{\bot1} \step s'_1$ and $s_2 = v_{\bot2} \step s'_2$ be two streams such
              that $s_1 \equiv_{n+1} s_2$. \Cref{def:semantics:grsync:theorems:determinism:equivn} ensures
              that $v_{\bot1} = v_{\bot2}$ and $s'_1 \equiv_n s'_2$. The induction hypothesis
              implies $s'_2 \equiv_n s'_1$. Because the equality on values is symmetrical, we can
              reuse \Cref{def:semantics:grsync:theorems:determinism:equivn} to state $s_2 \equiv_{n+1} s_1$.

        \item Transitive: \Cref{def:semantics:grsync:theorems:determinism:equiv0} implies for all streams
              $s_1$, $s_2$ and $s_3$ we have $s_1 \equiv_0 s_2$, $s_2 \equiv_0 s_3$ and
              $s_1 \equiv_0 s_3$.

              Let $n \in \N$ and suppose that for all streams $s_1$, $s_2$ and $s_3$ such
              that $s_1 \equiv_n s_2$ and $s_2 \equiv_n s_3$ we have $s_1 \equiv_n s_3$.
              Let $s_1 = v_{\bot1} \step s'_1$, $s_2 = v_{\bot2} \step s'_2$ and
              $s_3 = v_{\bot3} \step s'_3$ be streams such that $s_1 \equiv_{n+1} s_2$ and
              $s_2 \equiv_{n+1} s_3$. \Cref{def:semantics:grsync:theorems:determinism:equivn} ensures
              that $v_{\bot1} = v_{\bot2}$, $v_{\bot2} = v_{\bot3}$, $s'_1 \equiv_n s'_2$, and
              $s'_2 \equiv_n s'_3$. The induction hypothesis implies $s'_1 \equiv_n s'_3$. Because
              the equality on values is transitive, we can reuse
              \Cref{def:semantics:grsync:theorems:determinism:equivn} to state $s_1 \equiv_{n+1} s_3$.
    \end{itemize}
\end{skippableproof}

The proof of component determinism (\Cref{thm:semantics:grsync:theorems:determinism}) proceeds by
induction on $n$, the length of the stream prefixes being compared. The inductive step, in turn
requires a structural induction on the syntax of \GRSync. The
intermediate \Crefrange{lem:semantics:grsync:theorems:determinism:e}
{lem:semantics:grsync:theorems:determinism:eq} establish equivalent determinism properties
for each syntactic construct (expressions, patterns,
event patterns, and equations).
%
This nested proof structure leads to a mutual recursion, similar to the proof of productivity
(\Cref{sec:semantics:grsync:theorems:productivity}). The mutual recursion arises because the
structural induction on an expression may encounter a component call, which requires invoking the
main determinism theorem itself. This recursion is well-founded, however, because \GRSync
components are not mutually recursive, which guarantees that the inductive argument terminates.

\begin{lemma}[Determinism of \GRSync expressions]
    \label{lem:semantics:grsync:theorems:determinism:e}
    Let $e$ be a well-typed expression in a well-typed and causal program \G
    and a typing context \tyCtx. If $H_1$ and $H_2$ are stream contexts on clock $ck$, then for
    all $n \in \N$, and lists of streams $s_1^+$ and $s_2^+$:
    \begin{equation*}
        \wedge\begin{array}{l}
            \exprDenot{H_1}{ck}{e}{s_1^+} \\
            \exprDenot{H_2}{ck}{e}{s_2^+}
        \end{array}
        \; \Rightarrow \;
        \left(\forall y \in \deps{e}, \, H_1[y] \equiv_n H_2[y]\right)
        \; \Rightarrow \;
        (s_1 \equiv_n s_2)^+.
    \end{equation*}
\end{lemma}
\begin{skippableproof}
    Let $n \in \N$.
    The proof is done by structural induction on \GRSync expressions, and by
    induction on components (recursive component definitions are not allowed).

    \paragraph{Constants}
    Let $c$ be a constant expression. For all stream contexts $H_1$ and $H_2$ on clock $ck$ and
    list of streams $s_1^+$ and $s_2^+$ such that:
    \begin{align*}
         & \exprDenot{H_1}{ck}{c}{s_1^+}  \\
         & \exprDenot{H_2}{ck}{c}{s_2^+}.
    \end{align*}
    The rule \nameref{rule:semantics:grsync:denot:e:const} implies $[s_1^+] \equiv
        [\const{c}{ck}]$ and $[s_2^+] \equiv [\const{c}{ck}]$. By symmetry and
    transitivity we have $(s_1 \equiv_n s_2)^+$.

    \paragraph{Identifiers}

    Let $x$ be an identifier expression. The only dependency of the expression $x$ is the identifier
    $x$ (\Cref{rule:semantics:grsync:causality:dependencies:deps:e:ident}).
    %
    Let $H_1$ and $H_2$ denote stream contexts on clock $ck$ such that $H_1[x] \equiv_n H_2[x]$, and
    list of streams $s_1^+$ and $s_2^+$ such that:
    \begin{align*}
         & \exprDenot{H_1}{ck}{x}{s_1^+}  \\
         & \exprDenot{H_2}{ck}{x}{s_2^+}.
    \end{align*}
    The rule \nameref{rule:semantics:grsync:denot:e:ident} implies $[s_1^+] \equiv [H_1[x]]$ and
    $[s_2^+] \equiv [H_2[x]]$. By symmetry and transitivity we have $(s_1 \equiv_n s_2)^+$.

    \paragraph{Last memories}

    Let \last{x} be the memory of identifier $x$. The only dependency of the expression \last{x} is
    the identifier \last{x} (\Cref{rule:semantics:grsync:causality:dependencies:deps:e:last}).
    %
    Let $H_1$ and $H_2$ denote stream contexts on clock $ck$ such that $H_1[\last{x}] \equiv_n
        H_2[\last{x}]$, and list of streams $s_1^+$ and $s_2^+$ such that:
    \begin{align*}
         & \exprDenot{H_1}{ck}{\last{x}}{s_1^+}  \\
         & \exprDenot{H_2}{ck}{\last{x}}{s_2^+}.
    \end{align*}
    The rule \nameref{rule:semantics:grsync:denot:e:last} implies $[s_1^+] \equiv [H_1[\last{x}]]$
    and $[s_2^+] \equiv [H_2[\last{x}]]$. By symmetry and transitivity we have
    $(s_1 \equiv_n s_2)^+$.

    \paragraph{Event emission}

    Let \emit{e} be the emission of the well-typed expression $e$. Let $\{y^+\} = \deps{\emit{e}}$
    the set of dependencies of \emit{e}. Let $H_1$ and $H_2$ denote stream contexts on clock $ck$
    such that $(H_1[y] \equiv_n H_2[y])^+$, and list of streams $s_1^+$ and $s_2^+$ such that:
    \begin{align*}
         & \exprDenot{H_1}{ck}{\emit{e}}{s_1^+}  \\
         & \exprDenot{H_2}{ck}{\emit{e}}{s_2^+}.
    \end{align*}
    Because $\deps{\emit{e}} = \deps{e}$ (\Cref{rule:semantics:grsync:causality:dependencies:deps:e%
        :emit}), the structural induction hypothesis applied to $e$ with $H_1$ and $H_2$ on $ck$
    defines streams $s'_1$ and $s'_2$ such that:
    \begin{align*}
         & \exprDenot{H_1}{ck}{e}{s_1^+} \\
         & \exprDenot{H_2}{ck}{e}{s_2^+} \\
         & s'_1 \equiv_n s'_2.
    \end{align*}
    The rule \nameref{rule:semantics:grsync:denot:e:emit} implies $[s_1^+] \equiv
        [\some{s'_1}]$ and $[s_2^+] \equiv [\some{s'_2}]$, and we have $\some{s'_1}
        \equiv \some{s'_2}$. By symmetry and transitivity we have $(s_1 \equiv_n s_2)^+$.

    \paragraph{Operations}

    Let $op(e^+)$ be the well-typed application of operation $op$ on the well-typed expressions
    $e^+$. Let $\{y^+\} = \deps{op(e^+)}$ the set of dependencies of $op(e^+)$. Let $H_1$ and $H_2$
    denote stream contexts on clock $ck$ such that $(H_1[y] \equiv_n H_2[y])^+$, and list of streams
    $s_1^+$ and $s_2^+$ such that:
    \begin{align*}
         & \exprDenot{H_1}{ck}{op(e^+)}{s_1^+}  \\
         & \exprDenot{H_2}{ck}{op(e^+)}{s_2^+}.
    \end{align*}
    Because $\deps{op(e^+)} = \bigcup \deps{e}^+$ (\Cref{rule:semantics:grsync:causality:%
        dependencies:deps:e:op}), for all expression $e$ of $e^+$ and all $z \in \deps{e}$, we have
    $H_1[z] \equiv_n H_2[z]$. The structural induction hypothesis applied to $e^+$ defines streams
    $s^{\prime+}_1$ and $s^{\prime+}_2$ such that:
    \begin{align*}
         & (\exprDenot{H_1}{ck}{e}{s'_1})^+ \\
         & (\exprDenot{H_2}{ck}{e}{s'_2})^+ \\
         & (s'_1 \equiv_n s'_2)^+.
    \end{align*}
    The rule \nameref{rule:semantics:grsync:denot:e:op} implies $[s_1^+] \equiv [op(s^{\prime+}_1)]$
    and $[s_2^+] \equiv [op(s^{\prime+}_2)]$, and we have $op(s^{\prime+}_1) \equiv
        op(s^{\prime+}_2)$. By symmetry and transitivity we have $(s_1 \equiv_n s_2)^+$.

    \paragraph{Component calls}

    Let $f(e^+)$ be the call of the well-typed and causal component $f$ with the well-typed
    expressions $e^+$. Let $\{y^+\} = \deps{f(e^+)}$ the set of dependencies of $f(e^+)$. Let $H_1$
    and $H_2$ denote stream contexts on clock $ck$ such that $(H_1[y] \equiv_n H_2[y])^+$, and list
    of streams $s_1^+$ and $s_2^+$ such that:
    \begin{align*}
         & \exprDenot{H_1}{ck}{f(e^+)}{s_1^+}  \\
         & \exprDenot{H_2}{ck}{f(e^+)}{s_2^+}.
    \end{align*}
    Because $\deps{f(e^+)} = \bigcup \deps{e}^+$ (\Cref{rule:semantics:grsync:causality:%
        dependencies:deps:e:app}), for all expression $e$ of $e^+$ and all $z \in \deps{e}$, we have
    $H_1[z] \equiv_n H_2[z]$. The structural induction hypothesis applied to $e^+$ defines streams
    $s^{\prime+}_1$ and $s^{\prime+}_2$ such that:
    \begin{align*}
         & (\exprDenot{H_1}{ck}{e}{s'_1})^+ \\
         & (\exprDenot{H_2}{ck}{e}{s'_2})^+ \\
         & (s'_1 \equiv_n s'_2)^+.
    \end{align*}
    The induction on component calls, using the rule \nameref{rule:semantics:grsync:denot:e:app} and
    \Cref{thm:semantics:grsync:theorems:determinism}, states:
    $$
        \wedge \begin{array}{l}
            \compDenot{ck}{f}{s_1^{\prime+}}{s_{1}^+} \\
            \compDenot{ck}{f}{s_2^{\prime+}}{s_{2}^+}
        \end{array}
        \; \Rightarrow \;
        \left[(s'_1 \equiv_n s'_2)^+ \Rightarrow (s_1 \equiv_n s_2)^+\right],
    $$
    which implies $(s_1 \equiv_n s_2)^+$.

    \setcounter{paragraph}{0}
\end{skippableproof}

\begin{lemma}[Determinism of \GRSync patterns]
    \label{lem:semantics:grsync:theorems:determinism:pat}
    Let \guarded{pat^+ = s^+}{e_g} be a well-typed guarded pattern matching the streams $s^+$ in a
    program \G and a typing context \tyCtx. If $H_1$ and $H_2$ are stream contexts on clock $ck$,
    then for all $n \in \N$, and clocks $ck_1$ and $ck_2$:
    \begin{multline*}
       \patDenot{H_1}{ck}{\guarded{pat^+ = s^+}{e_g}}{ck_1} \quad \wedge \quad  \patDenot{H_2}{ck}{\guarded{pat^+ = s^+}{e_g}}{ck_2} \\
        \qquad \Rightarrow \;
        \left[\forall y \in \deps{e_g} \backslash \defs{pat^+}, \, H_1[y] \equiv_n H_2[y]\right] \hfill
        \\ \Rightarrow \;
        \left[
            ck_1 \equiv_n ck_2 \, \wedge \, \forall x \in \defs{pat^+}, \, \whenT{H_1}{ck_1}[x]
            \equiv_n \whenT{H_2}{ck_2}[x]
            \right].
    \end{multline*}
\end{lemma}
\begin{skippableproof}
    Let $n \in \N$. The proof is done by structural induction on \GRSync patterns, and
    relies on \Cref{lem:semantics:grsync:theorems:determinism:e} that states the determinism of \GRSync
    expressions.

    \paragraph{Guarded patterns}

    Let $\guarded{pat^+ = s^+}{e_g}$ be a well-typed pattern matching the streams $s^+$ and guarded
    by the expression $e_g$ in a program \G and a typing context \tyCtx. The pattern
    being well-typed, it ensures that each individual pattern $pat$ corresponds to a stream $s$, and
    that the expression $e_g$ is a boolean. The second point implies that any stream associated with
    $e_g$ is a clock.
    %
    Let $H_1$ and $H_2$ denote stream contexts on $ck$, and let clocks $ck_1$ and $ck_2$ such that:
    \begin{align*}
         & \patDenot{H_1}{ck}{\guarded{pat^+ = s^+}{e_g}}{ck_1}                         \\
         & \patDenot{H_2}{ck}{\guarded{pat^+ = s^+}{e_g}}{ck_2}                         \\
         & \forall y \in \deps{e_g} \backslash \defs{pat^+}, \, H_1[y] \equiv_n H_2[y].
    \end{align*}
    %
    The rule \nameref{rule:semantics:grsync:denot:pats} implies there exist clocks $ck^{\prime+}_1$,
    $ck^{\prime+}_2$, $ck^{pat}_1$, and $ck^{pat}_2$ such that:
    \begin{align*}
         & \left(\patDenot{H}{ck}{pat = s}{ck'_1}\right)^+                                            \\
         & \left(\patDenot{H}{ck}{pat = s}{ck'_2}\right)^+                                            \\
         & ck^{pat}_1 \equiv \kwf{\&\&}(ck^{\prime+}_1)                                               \\
         & ck^{pat}_2 \equiv \kwf{\&\&}(ck^{\prime+}_2)                                               \\
         & \exprDenot{\whenT{H_1}{ck^{pat}_1}}{\whenT{ck}{ck^{pat}_1}}{e_g}{\whenT{ck_1}{ck^{pat}_1}} \\
         & \exprDenot{\whenT{H_2}{ck^{pat}_2}}{\whenT{ck}{ck^{pat}_2}}{e_g}{\whenT{ck_2}{ck^{pat}_2}} \\
         & \whenF{ck_1}{ck^{pat}_1} \equiv \const{\false}{\whenF{ck}{ck^{pat}_1}}                     \\
         & \whenF{ck_2}{ck^{pat}_2} \equiv \const{\false}{\whenF{ck}{ck^{pat}_2}}
    \end{align*}
    %
    The induction hypothesis on patterns ensures $(ck'_1 \equiv_n ck'_2)^+$ which implies, by the
    determinism of functional applications on streams, that $\kwf{\&\&}(ck^{\prime+}_1) \equiv_n
        \kwf{\&\&}(ck^{\prime+}_2)$, thus $ck^{pat}_1 \equiv_n ck^{pat}_2$. The induction hypothesis
    also implies that for all pattern $pat$ and all $x$ in $\defs{pat}$, we have $\whenT{H_1}{ck'_1}
        [x] \equiv_n \whenT{H_2}{ck'_2}[x]$, so $\whenT{H_1}{ck^{pat}_1}[x]
        \equiv_n \whenT{H_2}{ck^{pat}_2}[x]$.
    %
    \Cref{rule:semantics:grsync:causality:defs:pat} ensures that $\defs{pat^+} = \bigcup
        \defs{pat}^+$, then for all $x$ in \defs{pat^+} we have $\whenT{H_1}{ck^{pat}_1}[x]
        \equiv_n \whenT{H_2}{ck^{pat}_2}[x]$. Therefore, for all $y \in \deps{e_g}$, we have
    $\whenT{H_1}{ck^{pat}_1}[y] \equiv_n \whenT{H_2}{ck^{pat}_2}[y]$.
    %
    \Cref{lem:semantics:grsync:theorems:determinism:e} with the contexts \whenT{H_{1,2}}{ck^{pat}_{1,2}} and
    with the clocks $\whenT{ck}{ck^{pat}_{1,2}}$ states:
    \begin{equation*}
        \whenT{ck_1}{ck^{pat}_1} \equiv_n \whenT{ck_2}{ck^{pat}_2}.
    \end{equation*}
    %
    Because $ck^{pat}_1 \equiv_n ck^{pat}_2$, the constant streams \const{\false}{\whenF{ck}
        {ck^{pat}_{1,2}}} are equivalent up-to-$n$. By symmetry and transitivity
    we have: $$\whenF{ck_1}{ck^{pat}_1} \equiv_n \whenF{ck_2}{ck^{pat}_2}.$$
    %
    The merge of both filtering on $ck$ implies $ck_1 \equiv_n ck_2$.

    \paragraph{Constants}

    Let $c = s$ be a constant pattern matching the stream $s$. It has no dependency.
    %
    Let $H_1$ and $H_2$ denote stream contexts on $ck$, and let clocks $ck_1$ and $ck_2$ such that:
    \begin{align*}
         & \patDenot{H_1}{ck}{c = s}{ck_1}  \\
         & \patDenot{H_2}{ck}{c = s}{ck_2}.
    \end{align*}
    %
    The rule \nameref{rule:semantics:grsync:denot:pat:const} implies $ck_1 \equiv
        (\cdot \kwf{==} c)(s)$ and $ck_2 \equiv (\cdot \kwf{==} c)(s)$. \Cref{rule:semantics:grsync%
        :causality:defs:pat:const} ensures that $\defs{c} = \emptyset$. By symmetry and transitivity
    we have $ck_1 \equiv_n ck_2$.

    \paragraph{Identifiers}

    Let $x = s$ be an identifier pattern matching the stream $s$. It has no dependency.
    %
    Let $H_1$ and $H_2$ denote stream contexts on $ck$, and let clocks $ck_1$ and $ck_2$ such that:
    \begin{align*}
         & \patDenot{H_1}{ck}{x = s}{ck_1}  \\
         & \patDenot{H_2}{ck}{x = s}{ck_2}.
    \end{align*}
    %
    The rule \nameref{rule:semantics:grsync:denot:pat:ident} implies $ck_1 \equiv\const{\true}{ck}$,
    $ck_2 \equiv \const{\true}{ck}$, $H_1[x] \equiv s$, and
    $H_2[x] \equiv s$. \Cref{rule:semantics:grsync:causality:defs:pat:ident} ensures that $\defs{x}
        = \{x\}$. By symmetry and transitivity we have $ck_1 \equiv_n ck_2$ and $H_1[x] \equiv_n
        H_2[x]$, which implies $\whenT{H_1}{ck_1}[x] \equiv_n \whenT{H_2}{ck_2}[x]$.

    \paragraph{Wildcards}

    Let $\_ = s$ be a wildcard pattern matching the stream $s$. It has no dependency.
    %
    Let $H_1$ and $H_2$ denote stream contexts on $ck$, and let clocks $ck_1$ and $ck_2$ such that:
    \begin{align*}
         & \patDenot{H_1}{ck}{\_ = s}{ck_1}  \\
         & \patDenot{H_2}{ck}{\_ = s}{ck_2}.
    \end{align*}
    %
    The rule \nameref{rule:semantics:grsync:denot:pat:wild} implies $ck_1 \equiv \const{\true}{ck}$
    and $ck_2 \equiv \const{\true}{ck}$. \Cref{rule:semantics:%
        grsync:causality:defs:pat:ident} ensures that $\defs{\_} = \emptyset$. By symmetry and
    transitivity we have $ck_1 \equiv_n ck_2$.

    \setcounter{paragraph}{0}
\end{skippableproof}

\begin{lemma}[Determinism of \GRSync event patterns]
    \label{lem:semantics:grsync:theorems:determinism:epat}
    Let \guarded{epat^+}{e_g} be a well-typed guarded event pattern in a program \G and a typing
    context \tyCtx. If $H_1$ and $H_2$ are stream contexts on clock $ck$, then for all
    $n \in \N$, and clocks $ck_1$ and $ck_2$:
    \begin{multline*}
        \wedge \begin{array}{l}
            \epatDenot{H_1}{ck}{\guarded{epat^+}{e_g}}{ck_1} \\
            \epatDenot{H_2}{ck}{\guarded{epat^+}{e_g}}{ck_2}
        \end{array}
        \; \Rightarrow \\
        \left[
        \forall y \in \deps{epat^+} \cup (\deps{e_g} \backslash \defs{epat^+}), \, H_1[y] \equiv_n H_2[y]
        \right]
        \; \Rightarrow \\
        \left[
            ck_1 \equiv_n ck_2 \, \wedge \, \forall x \in \defs{epat^+}, \, \whenT{H_1}{ck_1}
            [x] \equiv_n \whenT{H_2}{ck_2}[x]
            \right].
    \end{multline*}
\end{lemma}
\begin{skippableproof}
    Let $n \in \N$. The proof is done by structural induction on \GRSync event patterns, and
    relies on \Cref{lem:semantics:grsync:theorems:determinism:e} that states the determinism of \GRSync
    expressions.

    \paragraph{Guarded event patterns}

    Let \guarded{epat^+}{e_g} be a well-typed event pattern guarded by the expression $e_g$ in a
    program \G and a typing context \tyCtx. The event pattern being well-typed, it ensures that
    the expression $e_g$ is a boolean, thus any stream associated with $e_g$ is a clock.
    %
    Let $H_1$ and $H_2$ denote stream contexts on $ck$, and let clocks $ck_1$ and $ck_2$ such that:
    \begin{align*}
         & \epatDenot{H_1}{ck}{\guarded{epat^+}{e_g}}{ck_1}                                                   \\
         & \epatDenot{H_2}{ck}{\guarded{epat^+}{e_g}}{ck_2}                                                   \\
         & \forall y \in \deps{epat^+} \cup (\deps{e_g} \backslash \defs{epat^+}), \, H_1[y] \equiv_n H_2[y].
    \end{align*}
    %
    The rule \nameref{rule:semantics:grsync:denot:epats} implies there exist clocks $ck^{\prime+}_1$,
    $ck^{\prime+}_2$, $ck^{epat}_1$, and $ck^{epat}_2$ such that:
    \begin{align*}
         & \left(\patDenot{H}{ck}{epat}{ck'_1}\right)^+                                                  \\
         & \left(\patDenot{H}{ck}{epat}{ck'_2}\right)^+                                                  \\
         & ck^{epat}_1 \equiv \kwf{\&\&}(ck^{\prime+}_1)                                                 \\
         & ck^{epat}_2 \equiv \kwf{\&\&}(ck^{\prime+}_2)                                                 \\
         & \exprDenot{\whenT{H_1}{ck^{epat}_1}}{\whenT{ck}{ck^{epat}_1}}{e_g}{\whenT{ck_1}{ck^{epat}_1}} \\
         & \exprDenot{\whenT{H_2}{ck^{epat}_2}}{\whenT{ck}{ck^{epat}_2}}{e_g}{\whenT{ck_2}{ck^{epat}_2}} \\
         & \whenF{ck_1}{ck^{epat}_1} \equiv \const{\false}{\whenF{ck}{ck^{epat}_1}}                      \\
         & \whenF{ck_2}{ck^{epat}_2} \equiv \const{\false}{\whenF{ck}{ck^{epat}_2}}
    \end{align*}
    %
    \Cref{rule:semantics:grsync:causality:dependencies:deps:epat} states that $\deps{epat^+} =
        \bigcup \deps{epat}^+$, so we can apply the induction hypothesis on event patterns, ensuring
    $(ck'_1 \equiv_n ck'_2)^+$ and, by the determinism of functional applications on streams,
    $\kwf{\&\&}(ck^{\prime+}_1) \equiv_n \kwf{\&\&}(ck^{\prime+}_2)$, thus $ck^{epat}_1 \equiv_n
        ck^{epat}_2$. The induction hypothesis on event patterns also implies that for all event
    pattern $epat$ and all $x$ in \defs{epat}, $\whenT{H_1}{ck'_1}[x] \equiv_n \whenT{H_2}{ck'_2}
        [x]$, so $\whenT{H_1}{ck^{epat}_1}[x] \equiv_n \whenT{H_2}{ck^{epat}_2}[x]$.
    %
    \Cref{rule:semantics:grsync:causality:defs:epat} ensures that $\defs{epat^+} = \bigcup
        \defs{epat}^+$, then for all $x$ in \defs{epat^+}, $\whenT{H_1}{ck^{epat}_1}[x]
        \equiv_n \whenT{H_2}{ck^{epat}_2}[x]$. Therefore, for all identifier $y$ in
    $\deps{e_g}$, we have $\whenT{H_1}{ck^{epat}_1}[y] \equiv_n \whenT{H_2}{ck^{epat}_2}[y]$.
    We can apply \Cref{lem:semantics:grsync:theorems:determinism:e} with the contexts
    $\whenT{H_{1,2}}{ck^{epat}_{1,2}}$ and the clocks $\whenT{ck}{ck^{epat}_{1,2}}$ to
    state:
    \begin{equation*}
        \whenT{ck_1}{ck^{epat}_1} \equiv_n \whenT{ck_2}{ck^{epat}_2}.
    \end{equation*}
    %
    Because $ck^{epat}_1 \equiv_n ck^{epat}_2$, the constant streams \const{\false}{\whenF{ck}
        {ck^{epat}_{1,2}}} are equivalent up-to-$n$. By symmetry and transitivity
    we have: $$\whenF{ck_1}{ck^{epat}_1} \equiv_n \whenF{ck_2}{ck^{epat}_2}.$$
    %
    The merge of both filtering on $ck$ implies $ck_1 \equiv_n ck_2$.

    \paragraph{Event detection}

    Let $\detect{x}{y}$ be an event detection pattern.
    %
    Let $H_1$ and $H_2$ denote stream contexts on $ck$, and let clocks $ck_1$ and $ck_2$ such that:
    \begin{align*}
         & \epatDenot{H_1}{ck}{\detect{x}{y}}{ck_1}                       \\
         & \epatDenot{H_2}{ck}{\detect{x}{y}}{ck_2}                       \\
         & \forall y \in \deps{\detect{x}{y}}, \, H_1[y] \equiv_n H_2[y].
    \end{align*}
    %
    \Cref{rule:semantics:grsync:causality:dependencies:deps:epat:event} ensures
    $\deps{\detect{x}{y}}= \{y\}$, so we have $H_1[y] \equiv_n H_2[y]$.
    %
    Rule \nameref{rule:semantics:grsync:denot:epat:event} implies $ck_{1,2} \equiv
        \isSome{H_{1,2}[y]}$, and $\whenT{H_{1,2}[y]}{ck_{1,2}} \equiv
        \whenT{\some{H_{1,2}[x]}}{ck_{1,2}}$.
    %
    By symmetry and transitivity of the equivalence up-to-$n$, we have the equivalence of the
    associated clocks $ck_1 \equiv_n ck_2$, and $\whenT{\some{H_1[x]}}{ck_1} \equiv_n
        \whenT{\some{H_2[x]}}{ck_2}$.
    %
    \Cref{rule:semantics:grsync:causality:defs:epat:event} ensures that $\defs{\detect{x}{y}} =
        \{x\}$.

    \paragraph{Rising edges}

    Let $e$ be a rising edge pattern.
    %
    Let $H_1$ and $H_2$ denote stream contexts on $ck$, and let clocks $ck_1$ and $ck_2$ such that:
    \begin{align*}
         & \epatDenot{H_1}{ck}{e}{ck_1}                       \\
         & \epatDenot{H_2}{ck}{e}{ck_2}                       \\
         & \forall y \in \deps{e}, \, H_1[y] \equiv_n H_2[y].
    \end{align*}
    %
    The rule \nameref{rule:semantics:grsync:denot:epat:edge} implies there exist $ck'_1$ and $ck'_2$
    such that:
    \begin{align*}
         & \exprDenot{H_1}{ck}{e}{ck'_1}                               \\
         & \exprDenot{H_2}{ck}{e}{ck'_2}                               \\
         & ck_1 \equiv ck'_1 \kwf{\&\&} \kwf{!} \buffer{\false}{ck'_1} \\
         & ck_2 \equiv ck'_2 \kwf{\&\&} \kwf{!} \buffer{\false}{ck'_2}
    \end{align*}
    %
    \Cref{rule:semantics:grsync:causality:dependencies:deps:epat:edge} ensures
    $\deps{e}= \deps{e}$ (the dependencies of the expression), so we can apply \Cref{lem:semantics:grsync:%
        theorems:determinism:e} that states $ck'_1 \equiv_n ck'_2$.
    %
    By symmetry and transitivity we have $ck_1 \equiv_n ck_2$.
    %
    \Cref{rule:semantics:grsync:causality:defs:epat:edge} ensures that $\defs{\detect{x}{y}} =
        \emptyset$.

    \setcounter{paragraph}{0}
\end{skippableproof}

\begin{lemma}[Determinism of \GRSync equations]
    \label{lem:semantics:grsync:theorems:determinism:eq}
    Let $eq^+$ be a well-typed equation in a well-typed and causal program \G and a typing context
    \tyCtx. If $H_1$ and $H_2$ are stream contexts on clock $ck$, then for all $n \in \N$:
    \begin{multline*}
        \eqDenot{H_1}{ck}{eq^+} \, \wedge \, \eqDenot{H_2}{ck}{eq^+}
        \; \Rightarrow \\
        \left[\forall y \in \deps{eq^+}, \, H_1[y] \equiv_n H_2[y]\right]
        \; \Rightarrow \\
        \left[\forall x \in \defs{eq^+}, \, H_1[x] \equiv_n H_2[x]\right].
    \end{multline*}
\end{lemma}
\begin{skippableproof}
    Let $n \in \N$. The proof is done by structural induction on \GRSync equations, and
    relies on \Cref{lem:semantics:grsync:theorems:determinism:e,lem:semantics:grsync:theorems:determinism:pat,lem:%
        semantics:grsync:theorems:determinism:epat} that states the determinism of \GRSync expressions,
    patterns, and event patterns.

    \paragraph{Overview of the proof strategy}
    The difficult points of the proof are the following.
    \begin{itemize}
        \item For \kwf{match} equations, the exhaustiveness of patterns (verified by the \Rust
              compiler on the generated code) ensures that the clock filtering induced by each
              pattern matching is merged so that the entire match equation remains on the base clock.
        \item For \kwf{when} equations, it is a similar clock-branch filtering, but on which the
              stream of signals not updated in a branch is equivalent to their \kwf{last} stream, and
              the stream of events not emitted are stream repeating \none.
        \item For blocks of mutually recursive equations, we use rank induction
              (\Cref{def:semantics:grsync:causality:induction:rank}) to prove iteratively the
              equivalence of identifiers following the rank (\Cref{def:semantics:grsync:causality:rank}).
    \end{itemize}


    \paragraph{Declarations}
    Let $x^+ = e$ be a well-typed declaration in a well-typed and causal program \G and a typing
    context \tyCtx. Let $H_1$ and $H_2$ be stream contexts on a clock $ck$ such that:
    \begin{align*}
         & \eqDenot{H_1}{ck}{x^+ = e}                               \\
         & \eqDenot{H_2}{ck}{x^+ = e}                               \\
         & \forall y \in \deps{x^+ = e}, \, H_1[y] \equiv_n H_2[y].
    \end{align*}

    \nameref{rule:semantics:grsync:denot:eq:decl} ensures there exist streams $s_1^+$ and $s_2^+$
    such that:
    \begin{align*}
         & \exprDenot{H_1}{ck}{e}{s_1^+}  \\
         & \exprDenot{H_2}{ck}{e}{s_2^+}  \\
         & \left(H[x] \equiv s_1\right)^+ \\
         & \left(H[x] \equiv s_2\right)^+
    \end{align*}
    The declaration being well-typed, it ensures that for every $x$ corresponds a stream $s_1$ and
    $s_2$.

    \Cref{rule:semantics:grsync:causality:dependencies:deps:eq:decl} ensures that $\deps{x^+ = e} =
        \deps{e}$ so that we can apply \Cref{lem:semantics:grsync:theorems:determinism:e} on $e$ to state that
    we have $(s_1 \equiv_n s_2)^+$.

    \Cref{rule:semantics:grsync:causality:defs:eq:decl} ensures that $\defs{x^+ = e} = \{x^+\}$, and
    by symmetry and transitivity we have $(H_1[x] \equiv_n H_2[x])^+$.

    \paragraph{Match}
    Let \match{e}{x^+}{(\branch{pat_i^+}{e_i}{eq_i^+})^+} be a well-typed \kwf{match} equation in a
    well-typed and causal program \G under a typing context \tyCtx. Let $H_1$ and $H_2$ be stream
    contexts over a clock $ck$ such that:
    \begin{align*}
         & \eqDenot{H_1}{ck}{\match{e}{x^+}{(\branch{pat_i^+}{e_i}{eq_i^+})^+}}                               \\
         & \eqDenot{H_2}{ck}{\match{e}{x^+}{(\branch{pat_i^+}{e_i}{eq_i^+})^+}}                               \\
         & \forall y \in \deps{\match{e}{x^+}{(\branch{pat_i^+}{e_i}{eq_i^+})^+}}, \, H_1[y] \equiv_n H_2[y].
    \end{align*}

    \nameref{rule:semantics:grsync:denot:eq:match} ensures there exist streams $s_1^+$ and $s_2^+$,
    clocks $ck_{i, 1}^+$, $ck_{i, 2}^+$, $ck_{i, 1}^{\prime+}$ and $ck_{i, 2}^{\prime+}$, and stream
    contexts $H^\T_{i, 1}$ and $H^\T_{i, 2}$ such that:
    \begin{align*}
         & \exprDenot{H_1}{ck}{e}{s_1^+}                                                                \\
         & \exprDenot{H_2}{ck}{e}{s_2^+}                                                                \\
         & \left(\patDenot{H_1}{ck}{\guarded{pat_i^+ = s_1^+}{e_i}}{ck_{i, 1}}\right)^+                 \\
         & \left(\patDenot{H_2}{ck}{\guarded{pat_i^+ = s_2^+}{e_i}}{ck_{i, 2}}\right)^+                 \\
         & ck \equiv \kwf{\Vert}(ck_{i, 1}^+)                                                           \\
         & ck \equiv \kwf{\Vert}(ck_{i, 2}^+)                                                           \\
         & \left(ck'_{i, 1} \equiv ck_{i, 1} \kwf{\&\&} (\kwf{\&\&}(\kwf{!}ck_{j, 1})^{j < i})\right)^+ \\
         & \left(ck'_{i, 2} \equiv ck_{i, 2} \kwf{\&\&} (\kwf{\&\&}(\kwf{!}ck_{j, 2})^{j < i})\right)^+ \\
         & \left(H^\T_{i, 1} \equiv \whenT{H_1}{ck'_{i, 1}}\right)^+                                    \\
         & \left(H^\T_{i, 2} \equiv \whenT{H_2}{ck'_{i, 2}}\right)^+                                    \\
         & \left(\eqDenot{H^\T_{i, 1}}{ck'_{i, 1}}{eq_i^+}\right)^+                                     \\
         & \left(\eqDenot{H^\T_{i, 2}}{ck'_{i, 2}}{eq_i^+}\right)^+
    \end{align*}

    \Cref{rule:semantics:grsync:causality:dependencies:deps:eq:match} ensures that: $$\deps{\kwf
            {match}\{\dots\}} = \deps{e} \cup \bigcup \left(\deps{e_i} \cup \deps{eq_i^+} \backslash
        \defs{pat_i^+}\right)^+,$$ so that we can apply \Cref{lem:semantics:grsync:theorems:determinism:e} on
    $e$ to state that we have $(s_1 \equiv_n s_2)^+$.

    For all branch $i$ and all $y$ in $\deps{e_i} \backslash \defs{pat_i^+}$ we have $H_1[y]
        \equiv_n H_2[y]$, so \Cref{lem:semantics:grsync:theorems:determinism:pat} ensures that for all $i$:
    \begin{align*}
         & ck_{i, 1} \equiv_n ck_{i, 2}                                                                   \\
         & \forall x \in \defs{pat_i^+}, \, \whenT{H_1}{ck_{i, 1}}[x] \equiv_n \whenT{H_2}{ck_{i, 2}}[x].
    \end{align*}
    It implies that for all branch $i$, we have
    $ck_{i, 1} \kwf{\&\&} (\kwf{\&\&}(\kwf{!}ck_{j, 1})^{j < i}) \equiv_n
        ck_{i, 2} \kwf{\&\&} (\kwf{\&\&}(\kwf{!}ck_{j, 2})^{j < i})$ which ensures by symmetry and
    transitivity that $ck'_{i, 1} \equiv_n ck'_{i, 2}$, and that for all $x \in \defs{pat_i^+}$ we
    have $\whenT{H_1}{ck'_{i, 1}}[x] \equiv_n \whenT{H_2}{ck'_{i, 2}}[x]$ and thus
    $H^\T_{i, 1}[x] \equiv_n H^\T_{i, 2}[x]$.

    It follows that for all branch $i$ and $y$ in \deps{eq_i^+} we have $H^\T_{i, 1}[y] \equiv_n
        H^\T_{i, 2}[y]$, so we can apply the induction hypothesis on equations to state that for all
    $x$ in \defs{eq_i^+} we have $H^\T_{i, 1}[x] \equiv_n H^\T_{i, 2}[x]$.

    Since patterns are exhaustive, we have $ck \equiv \Vert(ck_{i, 1}^+)$ and $ck \equiv \Vert
        (ck_{i, 2}^+)$, meaning that for every time step, some pattern is active. And because this
    \kwf{match} equation is well-typed, all equation blocks defines the same identifiers. This
    ensures that for all $x$ in \defs{\kwf{match}\{\dots\}} we have $H_1[x] \equiv_n H_2[x]$.

    \paragraph{When}
    Let \when{x^+}{x_m^+ = c^+}{(\branch{epat_i^+}{e_i}{eq_i^+})^+} well-typed, in a\\well-typed and
    causal program \G under a typing context \tyCtx.

    Since the equations $eq_i^+$ may depend on memories of signals defined by the \kwf{when}
    equation $(\last{x_m})^+$, we prove by induction on $n$ that for all $x$ in \defs{\kwf{when}\{
        \dots \}} we have $H_1[x] \equiv_n H_2[x]$.

    \begin{itemize}
        \item \textit{Base case:} At $n = 0$, all pairs of streams are in $\equiv_0$, by definition
              of $\equiv_0$. In particular, we have $H_1[x] \equiv_0 H_2[x]$ for all contexts
              $H_1$ and $H_2$ and all $x$ in \defs{\kwf{when}\{ \dots \}}

        \item \textit{Inductive step:} Let $H_1$ and $H_2$ be stream contexts over
              a clock $ck$, such that:
              \begin{align*}
                   & \eqDenot{H_1}{ck}{\when{x^+}{x_m^+ = c^+}{(\branch{epat_i^+}{e_i}{eq_i^+})^+}} \\
                   & \eqDenot{H_2}{ck}{\when{x^+}{x_m^+ = c^+}{(\branch{epat_i^+}{e_i}{eq_i^+})^+}} \\
                   & \forall y \in \deps{\kwf{when}\{ \dots \}}, \, H_1[y] \equiv_{n+1} H_2[y].
              \end{align*}
              It implies:
              \begin{equation*}
                  \forall y \in \deps{\kwf{when}\{ \dots \}}, \, H_1[y] \equiv_n H_2[y],
              \end{equation*}
              so we suppose, as induction hypothesis, that:
              \begin{equation*}
                  \forall x \in \defs{\kwf{when}\{ \dots \}}, \, H_1[x] \equiv_n H_2[x].
              \end{equation*}

              \nameref{rule:semantics:grsync:denot:eq:when} ensures there exist clocks $ck_{i, 1}^+$,
              $ck_{i, 2}^+$, $ck_{i, 1}^{\prime+}$ and $ck_{i, 2}^{\prime+}$, and stream contexts
              $H^\T_{i, 1}$, $H^\T_{i, 2}$,$H^\F_1$ and $H^\F_2$ such that:
              \begin{align*}
                   & \left(\patDenot{H_1}{ck}{\guarded{epat_i^+ = s_1^+}{e_i}}{ck_{i, 1}}\right)^+                                                           \\
                   & \left(\patDenot{H_2}{ck}{\guarded{epat_i^+ = s_2^+}{e_i}}{ck_{i, 2}}\right)^+                                                           \\
                   & \left(ck'_{i, 1} \equiv ck_{i, 1} \kwf{\&\&} (\kwf{\&\&}(\kwf{!}ck_{j, 1})^{j < i})\right)^+                                            \\
                   & \left(ck'_{i, 2} \equiv ck_{i, 2} \kwf{\&\&} (\kwf{\&\&}(\kwf{!}ck_{j, 2})^{j < i})\right)^+                                            \\
                   & \left(H^\T_{i, 1} \equiv \whenT{H_1}{ck'_{i, 1}}\right)^+                                                                               \\
                   & \left(H^\T_{i, 2} \equiv \whenT{H_2}{ck'_{i, 2}}\right)^+                                                                               \\
                   & \left(\forall x_s \in \{x_m^+\} \backslash \defs{eq_i^+}, H^\T_{i, 1}[x_s] \equiv H^\T_{i, 1}[\last{x_s}]\right)^+                      \\
                   & \left(\forall x_s \in \{x_m^+\} \backslash \defs{eq_i^+}, H^\T_{i, 2}[x_s] \equiv H^\T_{i, 2}[\last{x_s}]\right)^+                      \\
                   & \left(\forall x_e \in \{x^+\} \backslash \{x_m^+\} \backslash \defs{eq_i^+}, H^\T_{i, 1}[x_e] \equiv \const{\none}{ck'_{i, 1}}\right)^+ \\
                   & \left(\forall x_e \in \{x^+\} \backslash \{x_m^+\} \backslash \defs{eq_i^+}, H^\T_{i, 2}[x_e] \equiv \const{\none}{ck'_{i, 2}}\right)^+ \\
                   & \left(\eqDenot{H^\T_{i, 1}}{ck'_{i, 1}}{eq_i^+}\right)^+                                                                                \\
                   & \left(\eqDenot{H^\T_{i, 2}}{ck'_{i, 2}}{eq_i^+}\right)^+                                                                                \\
                   & ck^\F_1 \equiv \whenF{ck}{\kwf{\Vert}(ck_{i, 1}^+)}                                                                                     \\
                   & ck^\F_2 \equiv \whenF{ck}{\kwf{\Vert}(ck_{i, 2}^+)}                                                                                     \\
                   & H^\F_1 \equiv \whenF{H}{\kwf{\Vert}(ck_{i, 1}^+)}                                                                                       \\
                   & H^\F_2 \equiv \whenF{H}{\kwf{\Vert}(ck_{i, 2}^+)}                                                                                       \\
                   & \forall x_s \in \{x_m^+\}, H^\F_1[x_s] \equiv H^\F_1[\last{x_s}]                                                                        \\
                   & \forall x_s \in \{x_m^+\}, H^\F_2[x_s] \equiv H^\F_2[\last{x_s}]                                                                        \\
                   & \forall x_e \in \{x^+\} \backslash \{x_m^+\}, H^\F_1[x_e] \equiv \const{\none}{ck^\F_1}                                                 \\
                   & \forall x_e \in \{x^+\} \backslash \{x_m^+\}, H^\F_2[x_e] \equiv \const{\none}{ck^\F_2}                                                 \\
                   & \left(H_1[\last{x_m}] \equiv \buffer{c}{H_1[x_m]} \right)^+                                                                             \\
                   & \left(H_2[\last{x_m}] \equiv \buffer{c}{H_2[x_m]} \right)^+
              \end{align*}

              \Cref{rule:semantics:grsync:causality:dependencies:deps:eq:when} ensures that $$\deps{
                      \kwf{when}\{\dots\}} = \begin{array}{l}
                      \bigcup \left(\deps{epat_i^+} \cup \deps{e_i} \cup \deps{eq_i^+} \backslash \defs{epat_i^+}\right)^+ \\
                      \qquad \backslash (\{x^+\} \cup \{\last{x_m}^+\})
                  \end{array}.$$

              For all branch $i$ and all $y$ in $\deps{e_i} \backslash \defs{epat_i^+}$ we have
              $H_1[y] \equiv_{n+1} H_2[y]$, so \Cref{lem:semantics:grsync:theorems:determinism:epat} ensures
              that for all $i$:
              \begin{align*}
                   & ck_{i, 1} \equiv_{n+1} ck_{i, 2}                                         \\
                   & \forall x \in \defs{epat_i^+}, \, \whenT{H_1}{ck_{i, 1}}[x] \equiv_{n+1}
                  \whenT{H_2}{ck_{i, 2}}[x].
              \end{align*}

              It implies that for all $i$, we have $ck_{i, 1} \kwf{\&\&} (\kwf{\&\&}(\kwf{!}ck_{j,
                      1})^{j < i}) \equiv_{n+1} ck_{i, 2} \kwf{\&\&} (\kwf{\&\&}(\kwf{!}ck_{j, 2}
                  )^{j < i})$ which ensures that $ck'_{i, 1} \equiv_{n+1} ck'_{i, 2}$, and that for
              all $x$ in $\defs{epat_i^+}$ we have
              $\whenT{H_1}{ck'_{i, 1}}[x]$ equivalent to $\whenT{H_2}{ck'_{i, 2}}[x]$
              up to $n+1$, and thus $H^\T_{i, 1}[x] \equiv_{n+1} H^\T_{i, 2}[x]$.

              Equivalently, we have $ck^\F_1 \equiv_{n+1} ck^\F_2$.

              The identifiers $x_m^+$ are included in
              \defs{\kwf{when}\{ \dots \}} (\Cref{rule:semantics:grsync:causality:defs:eq:when}).
              So for all $x_m$, we have $H_1[x_m] \equiv_n H_2[x_m]$ by the induction hypothesis on
              $n$. Since for each $j \in \{1,2\}$ we have $H_j[\last{x_m}] \equiv \buffer{c}
                  {H_j[x_m]}$, it follows that $H_1[\last{x_m}] \equiv_{n+1} H_2[\last{x_m}]$.

              For each $i$, for all $y$ in \deps{eq_i^+}, we have $H^\T_{i, 1}[y] \equiv_{n+1}
                  H^\T_{i, 2}[y]$. By induction hypothesis over equations, for all $x$ in
              \defs{eq_i^+}, we have $H^\T_{i, 1}[x] \equiv_{n+1} H^\T_{i, 2}[x]$.

              For each $i$, identifiers defined by the \kwf{when} equation and not updated by
              $eq_i^+$ ensure the following:
              \begin{itemize}
                  \item Signals $x_s$ are associated with their last values, so for each $j \in \{1,
                            2\}$ we have $H^\T_{i, j}[x_s] \equiv H^\T_{i, j}[\last{x_s}]$, and
                        because we proved the equivalence of \last{x_s} up to $n+1$ we can state:
                        $H^\T_{i, 1}[x_s] \equiv_{n+1} H^\T_{i, 2}[x_s]$.

                  \item Events $x_e$ are associated with a constant stream of \none values, so
                        for each $j \in \{1, 2\}$ we have $H^\T_{i, j}[x_e] \equiv \const{\none}
                            {ck'_{i, j}}$, and because we proved $ck'_{i, 1} \equiv_{n+1}
                            ck'_{i, 2}$ we have $H^\T_{i, 1}[x_e] \equiv_{n+1} H^\T_{i, 2}[x_e]$.
              \end{itemize}

              In the absence of event, the same rules apply:
              \begin{itemize}
                  \item Signals $x_s$ are associated with their last values, so for each $j \in \{1,
                            2\}$ we have $H^\F_j[x_s] \equiv H^\F_j[\last{x_s}]$, and because we
                        proved the equivalence of \last{x_s} up to $n+1$ we can state:  $H^\F_1[x_s]
                            \equiv_{n+1} H^\F_2[x_s]$.

                  \item Events $x_e$ are associated with a constant stream of \none values, so
                        for each $j \in \{1, 2\}$ we have $H^\F_j[x_e]\equiv\const{\none}{ck^\F_j}$,
                        and because we proved $ck^\F_1 \equiv_{n+1} ck^\F_2$ we have $H^\F_1[x_e]
                            \equiv_{n+1} H^\F_2[x_e]$.
              \end{itemize}

              Therefore, at each time step, whether an event pattern is triggered or not, the
              streams remain equivalent up to $n+1$: for all $x$ in \defs{\kwf{when} \{\dots\}} we
              have $H_1[x] \equiv_{n+1} H_2[x]$.
    \end{itemize}

    \paragraph{Multiple equations}
    Let $eq^+$ be a set of well-typed equations in a well-typed and causal program \G and a typing
    context \tyCtx. Let $H_1$ and $H_2$ be stream contexts over a clock $ck$ such that:
    \begin{align*}
         & \eqDenot{H_1}{ck}{eq^+}                               \\
         & \eqDenot{H_2}{ck}{eq^+}                               \\
         & \forall y \in \deps{eq^+}, \, H_1[y] \equiv_n H_2[y].
    \end{align*}

    The causality of $eq^+$ allows to use the rank induction
    (\Cref{def:semantics:grsync:causality:induction:rank}) to
    determine that for all $x \in \defs{eq^+} \cup \deps{eq^+}$, we have
    $H_1[x] \equiv_n H_2[x]$.

    \begin{itemize}
        \item \textit{Base case:} Let $x \in \defs{eq^+} \cup \deps{eq^+}$ such that
              $\rank{eq^+}{x} = 0$.

              Therefore, there exists an equation $eq_x$ that defines $x$ and
              that has no dependency with \defs{eq^+}. By induction hypothesis on equations
              structure, for all $z$ in \defs{eq_x} we have $H_1[z] \equiv_n H_2[z]$, in particular
              we have $H_1[x] \equiv_n H_2[x]$.

        \item \textit{Inductive step:} Let $k \in \N$.
              We assume that for all $y \in \defs{eq^+} \cup \deps{eq^+}$ such that
              $\rank{eq^+}{y} \le k$ we have $H_1[y] \equiv_n H_2[y]$.

              Let $x \in \defs{eq^+}$ such that $\rank{eq^+}{x} = k+1$.
              For all $y \in \defs{eq^+} \cup \deps{eq^+}$ such that $x$ depends on $y$ (\ie \depOrd{eq^+}{x}{y}),
              \Cref{lem:semantics:grsync:causality:rank:decrease}
              states that $\rank{eq^+}{x} > \rank{eq^+}{y}$, which implies that $\rank{eq^+}{y} \le k$.
              By rank induction hypothesis, for all $y \in \defs{eq^+} \cup \deps{eq^+}$ such that
              \depOrd{eq^+}{x}{y}, we have $H_1[y] \equiv_n H_2[y]$.

              Let $eq_x$ be the equation that defines $x$. Then for all $y \in \deps{eq_x}$, we have
              $H_1[y] \equiv_n H_2[y]$. By induction hypothesis on equations, for all $z$
              in \defs{eq_z} we have $H_1[z] \equiv_n H_2[z]$, in particular we have $H_1[x]
                  \equiv_n H_2[x]$.
    \end{itemize}

    \setcounter{paragraph}{0}
\end{skippableproof}

\begin{theorem}[Determinism]
    \label{thm:semantics:grsync:theorems:determinism}
    Let $f$ be a well-typed and causal component in a program \G and a typing
    context \tyCtx. If $[s_{i, 1}^+]$ and $[s_{i, 2}^+]$ are lists of input streams on clock $ck$,
    then for all $n \in \N$, lists of streams $[s_{o, 1}^+]$ and $[s_{o, 2}^+]$:
    $$
        \compDenot{ck}{f}{s_{i, 1}^+}{s_{o, 1}^+}
        \, \wedge \,
        \compDenot{ck}{f}{s_{i, 2}^+}{s_{o, 2}^+}
        \; \Rightarrow \;
        \left[
            \left(s_{i, 1} \equiv_n s_{i, 2}\right)^+
            \, \Rightarrow \,
            \left(s_{o, 1} \equiv_n s_{o, 2}\right)^+
            \right].
    $$
\end{theorem}
\begin{skippableproof}
    Let $f$ be a well-typed and causal component in a program \G and a typing context \tyCtx:
    $$\component{f}{x_i^+}{x_o^+}{eq^+}{x_m^+ = c^+}.$$
    Let $[s_{i, 1}^+]$, $[s_{i, 2}^+]$, $[s_{o, 1}^+]$, and $[s_{o, 2}^+]$ be lists of input streams
    on  clock $ck$ such that:
    \begin{align*}
         & \compDenot{ck}{f}{s_{i, 1}^+}{s_{o, 1}^+}  \\
         & \compDenot{ck}{f}{s_{i, 2}^+}{s_{o, 2}^+}.
    \end{align*}

    By application of the rule for components \nameref{rule:semantics:grsync:denot:comp}, it implies
    the existence of contexts $H_1$ and $H_2$ such that:
    \begin{align*}
         & \left(H_1[x_i] \equiv s_{i, 1}\right)^+                     \\
         & \left(H_2[x_i] \equiv s_{i, 2}\right)^+                     \\
         & \left(H_1[x_o] \equiv s_{o, 1}\right)^+                     \\
         & \left(H_2[x_o] \equiv s_{o, 2}\right)^+                     \\
         & \left(H_1[\last{x_m}] \equiv \buffer{c}{H_1[x_m]} \right)^+ \\
         & \left(H_2[\last{x_m}] \equiv \buffer{c}{H_2[x_m]} \right)^+ \\
         & \left(\eqDenot{H_1}{ck}{eq}\right)^+                        \\
         & \left(\eqDenot{H_2}{ck}{eq}\right)^+
    \end{align*}

    It is then sufficient to prove that, for all $n \in \N$, contexts $H_1$ and $H_2$ are
    equivalent up-to-$n$ on the domain of component's identifiers $\{x_i^+\} \cup \{\last{x_m}^+\}
        \cup \defs{eq^+}$:
    \begin{equation*}
        \P{n}:
        \left[\left(s_{i, 1} \equiv_n s_{i, 2}\right)^+ \Rightarrow [\forall x \in \{x_i^+\}
        \cup \{\last{x_m}^+\} \cup \defs{eq^+}, \, H_1[x] \equiv_n H_2[x]]\right].
    \end{equation*}
    The proof of \P{n} is done by induction on $n \in \N$ and relies on \Cref{lem%
        :semantics:grsync:theorems:determinism:eq} that states the determinism up-to-$n$ of \GRSync
    equations when their dependencies are determined.

    \begin{itemize}
        \item \textit{Base case:} At $n = 0$, all pairs of streams are in $\equiv_0$ (\Cref{def:%
                  semantics:grsync:theorems:determinism:equiv0}). In particular, for all $x$ in
              \defs{eq^+}, we have $H_1[x] \equiv_0 H_2[x]$. The property \P{0} holds.

        \item \textit{Inductive step:} Assuming \P{n} holds for $n \in \N$.
              Suppose $\left(s_{i, 1} \equiv_{n+1} s_{i, 2}\right)^+$. By definition of
              $\equiv_{n+1}$ (\Cref{def:semantics:grsync:theorems:determinism:equivn}), we have
              $\left(s_{i, 1} \equiv_n s_{i, 2}\right)^+$. The application of \P{n} with
              $[s_{i, 1}^+]$ and $[s_{i, 2}^+]$ states that, for all identifier $x$ in \defs{eq^+},
              we have $H_1[x] \equiv_n H_2[x]$.

              %% Inputs
              For $j \in \{1,2\}$, the streams associated with the inputs $(x_i^+)$ in $H_j$ are
              equivalent to $(s_{i, j}^+)$.
              We supposed $\left(s_{i, 1} \equiv_{n+1} s_{i, 2}\right)^+$, therefore for all input
              $x_i$, we have $H_1[x_i] \equiv_{n+1} H_2[x_i]$.

              %% Memories
              For $j \in \{1,2\}$, the stream associated with \last{x_m} in $H_j$ is equivalent to 
              the stream \buffer{c}{H_j[x_m]} (\cf rules \nameref{rule:semantics:grsync:denot:comp}
              and \nameref{rule:semantics:grsync:denot:eq:when}). The induction hypothesis
              \P{n} gives $H_1[x_m] \equiv_n H_2[x_m]$, which implies $\buffer{c}
                  {H_1[x_m]} \equiv_{n+1} \buffer{c}{H_2[x_m]}$. Therefore, for all memory
              $\last{x_m}$, we have $H_1[\last{x_m}] \equiv_{n+1} H_2[\last{x_m}]$.

              %% Others: by causality
              The causality of $f$ allows to use the rank induction
              (\Cref{def:semantics:grsync:causality:induction:rank}) to
              determine that for all $x \in \defs{eq^+} \cup \deps{eq^+}$, we have
              $H_1[x] \equiv_n H_2[x]$.

              \begin{itemize}
                  \item \textit{Base case:} Let $x \in \defs{eq^+} \cup \deps{eq^+}$ such that
                        $\rank{eq^+}{x} = 0$.

                        Therefore, there exists an equation $eq_x$ that defines $x$ and
                        that has no dependency with \defs{eq^+}.
                        \Cref{lem:semantics:grsync:theorems:determinism:eq} ensures that
                        for all $z$ in \defs{eq_x} we have $H_1[z] \equiv_n H_2[z]$, in particular
                        we have $H_1[x] \equiv_n H_2[x]$.

                  \item \textit{Inductive step:} Let $k \in \N$.
                        We assume that for all $y \in \defs{eq^+} \cup \deps{eq^+}$ such that
                        $\rank{eq^+}{y} \le k$ we have $H_1[y] \equiv_n H_2[y]$.

                        Let $x \in \defs{eq^+}$ such that $\rank{eq^+}{x} = k+1$.
                        For all $y \in \defs{eq^+} \cup \deps{eq^+}$ such that $x$ depends on $y$ (\ie \depOrd{eq^+}{x}{y}),
                        \Cref{lem:semantics:grsync:causality:rank:decrease}
                        states that $\rank{eq^+}{x} > \rank{eq^+}{y}$, which implies that $\rank{eq^+}{y} \le k$.
                        By rank induction hypothesis, for all $y \in \defs{eq^+} \cup \deps{eq^+}$ such that
                        \depOrd{eq^+}{x}{y}, we have $H_1[y] \equiv_n H_2[y]$.

                        Let $eq_x$ be the equation that defines $x$. Then for all $y \in \deps{eq_x}$, we have
                        $H_1[y] \equiv_n H_2[y]$. \Cref{lem:semantics:grsync:theorems:determinism:eq}
                        ensures that for all $z$ in \defs{eq_z} we have $H_1[z] \equiv_n H_2[z]$, in
                        particular we have $H_1[x] \equiv_n H_2[x]$.
              \end{itemize}

              %% Conclusion
              It follows that \P{{n+1}} holds.
    \end{itemize}

    Because $f$ is well-typed, the identifiers $\{x_o^+\}$ are included in \defs{eq^+}. Therefore,
    for all $n \in \N$ the property \P{n} implies that we have:
    \begin{equation*}
        \left(s_{i, 1} \equiv_n s_{i, 2}\right)^+ \Rightarrow \left(H_1[x_o] \equiv_n H_2[x_o]\right)^+
    \end{equation*}
    The rule for component semantics \nameref{rule:semantics:grsync:denot:comp} ensures for all
    $x_o$ that $H_1[x_o] \equiv s_{o, 1}$ and $H_2[x_o] \equiv s_{o, 2}$. By symmetry and
    transitivity, we have $\left(s_{o, 1} \equiv_n s_{o, 2}\right)^+$.
\end{skippableproof}

\paragraph{Summary}
In conclusion, this section formally established the determinism of \GRSync components. The main
result, presented in \Cref{thm:semantics:grsync:theorems:determinism}, proves that the output
streams of a component are uniquely determined by its input streams.
%
The proof was conducted using a coinductive argument, framed as an induction on the length of stream
prefixes. This approach was supported by a larger induction on the language constructs.
%
Similar to the productivity proof, the challenge of mutually recursive equations was addressed using
rank induction (\Cref{def:semantics:grsync:causality:induction:rank}), which leverages the causality
of the program to resolve dependencies in a well-defined order.
%
This guarantee of determinism is fundamental, ensuring that component behavior is predictable and
repeatable, which is a critical requirement for reliable reactive systems.
